\documentclass[10pt,a4paper]{amsart}
\usepackage[margin=19mm]{geometry}
\usepackage{amsmath,amssymb,mathtools}
\usepackage[scr=rsfs,cal=euler]{mathalfa}
\usepackage{xfrac}
\usepackage[unicode,hidelinks,bookmarks=false]{hyperref}
\newcommand{\ie}{i.e.\ }
\newcommand{\cf}{cf.\ }
\newcommand{\resp}{respectively\ }
\newcommand{\Subsection}{\subsection}
\providecommand{\nobreakdash}{\nobreak\hbox{-}}
\setlength{\emergencystretch}{2em}
\multlinegap0pt
\usepackage{amssymb}
\usepackage{tikz}
\newtheorem{proposition}{Proposition}
\newtheorem{theorem}{Theorem}
\def\bC{\mathbf{C}}
\def\bR{\mathbf{R}}
\newcommand\cT{\mathcal{T}}
\title{Legendrian and Lagrangian higher torsion}
\author{Daniel \'Alvarez-Gavela, Kiyoshi Igusa, Michael G Sullivan}
\date{Source: arXiv:2603.28007v2 (CC BY 4.0)}
\begin{document}
\maketitle

\noindent\emph{Source and licence.} Legendrian and Lagrangian higher torsion. Version arXiv:2603.28007v2 (CC BY 4.0), distributed by arXiv under the Creative Commons Attribution 4.0 International licence, http://creativecommons.org/licenses/by/4.0/, as recorded in the arXiv licence metadata for this exact version and shown on the abstract page https://arxiv.org/abs/2603.28007v2. Full licence notice: \texttt{licenses/arxiv-2603.28007-CC-BY-4.0.txt}.\par\smallskip

\noindent\textit{Editorial scope.} Counted unit: the article's proposition Leg (source0.tex lines 1261--1263). The article's complete original proof of it is delivered with it (source0.tex lines 1265--1273, one \texttt{proof} environment of 9 lines). No mathematics is written, repaired or re-derived: the statement and the proof are the article's own lines, in the article's own notation, and no retained line is substituted or re-flowed.  Retained uncounted as the source-local context the proof needs: lines 1150--1153; lines 1172--1176; lines 1252--1253; lines 5448--5453.  Nothing was omitted from any retained span, so the package carries no omission record.  No retained line contains a citation, so the wrapper carries no bibliography and no bibliography entry.  No retained line contains a figure environment, an image-inclusion command or drawing code, so no image is delivered and the package declares no asset dependency.  Editorial formatting only, in addition: an AMS \texttt{amsart} preamble that restates the article's own declarations and macros used by the retained lines, namely the environments \texttt{proposition}, \texttt{theorem} on separate counters, together with the standard packages those lines need.  The retained environments print in this document as Proposition 1, Theorem 1 in order of appearance, whereas the article prints the same items with the numbering of its own full document.  The complete native source of the article as distributed in the arXiv e-print for this exact version is bundled unchanged as \texttt{sources/197401/source0.tex}.
\begin{proposition}\label{prop: stab framing}
Let $f:M \times \bR^{4N} \to \bR$ be an orientable fibered function of tube type with moderate singularities, with associated fibered function of quadratic tube type $g:M \times \bR^{4N} \to \bR$ and tube bundle $W=T(g)$. The tube torsion $\tau(W) = \sum_k \tau_k(W)$ of $W$ may be computed from $f$ by the following formula: 
$$\tau(W) =    -\left(\tau_{2k}(f)  +   \pi_*^{\Sigma f} p_k(\gamma_f)\right).$$
\end{proposition}

\begin{proposition}\label{prop: effect twisted stab}
Let $f:M \times \bR^{4N_1} \to \bR$ be an orientable  fibered  function of tube type of even index and let $q:M \times \bR^{4N_2} \to \bR$ be an orientable fibered function of rigid tube type of even index. The higher torsions of the tube bundles $T(f)$ and $T(f \oplus q)$are related by the formula
$$ \tau_k(T(f \oplus q) ) = \tau_k(T(f))  - p_k(E) $$
where $E \to M$ is the stable bundle of $q$.
\end{proposition}

\begin{proposition}\label{prop: framing for special Leg}
Let $\Lambda \subset J^1M$ be a Legendrian whose front  $\Lambda \hookrightarrow J^1M \to J^0M$ only has cusp singularities and whose projection to the base $\Lambda \to M$ is a homotopy equivalence. If $\Lambda$ admits a generating function of tube type $f:M \times \bR^{4N_1} \to \bR$, then for a suitable fibered function of rigid tube type $q:M \times \bR^{4N_2} \to \bR$ the direct sum $f \oplus q : M \times \bR^{4N_1+4N_2} \to \bR$ is a frameable generating function of tube type , i.e. such that the stable bundle $\gamma_{f \oplus q}$ is trivial. \end{proposition}

\begin{theorem}\label{thm: framing} If  $f:W \to \bR$ is a function with moderate singularities, then
\begin{equation}\label{eq:FramingPrinciple}
\tau(W,\rho)=\sum_j\tau_j(f,\rho)+ \frac12 \sum_k (-1)^k \zeta(2k+1) \pi_\ast^{\Sigma f} \left( ch_{2k}(\gamma_f\otimes\bC)\right) 
\end{equation}
where the terms are defined below.
\end{theorem}

\begin{proposition}\label{prop: torsion for special Leg}
Let $\Lambda \subset J^1M$ be a Legendrian whose front projection $\Lambda \hookrightarrow J^1M \to J^0M$ only has cusp singularities and whose projection to the base $\Lambda \to M$ is a homotopy equivalence. The higher torsion of $\Lambda$ is the subset $\tau_{2k}(\Lambda) \subset \text{coker}(p_k)$ consisting of the cosets of $\text{im}(p_k)$ with representatives given by the cohomology classes $$\tau_{2k}(f)\in H^{4k}(M;\bR),$$ where $f:M \times \bR^{4N} \to \bR$ ranges through the collection of orientable generating functions of tube type (of even index) for $\Lambda$. 
\end{proposition}

\begin{proof}
First we show that for any such $f \times \bR^{4N} \to \bR$ we have $\tau_{2k}(f) \in \tau_{2k}(\Lambda)$. Now, we certainly have that $W=T(f)$ is in $\cT(\Lambda)$ and hence $-\tau_{2k}(W) \in \tau_{2k}(\Lambda)$. By Proposition \ref{prop: stab framing} we have
$$ -\tau_{2k}(W) =\tau_{2k}(f) + (\pi^{\Sigma f})_*p_k(\gamma_f). $$
Since $\pi^{\Sigma f}: \Sigma_f \to M$ is a homotopy equivalence, we have $\gamma_f=(\pi^{\Sigma f})^*E$ for some real vector bundle $E \to M$ and hence
$$ (\pi^{\Sigma f})_*p_k(\gamma_f) = (\pi^{\Sigma f})_*(\pi^{\Sigma f})^*p_k(E)=p_k(E), $$
where the last equality follows from the fact that $(\pi^{\Sigma f})_*(\pi^{\Sigma f})^*$ is multiplication by the Euler characteristic of the fiber, i.e. $+1$. Hence $\tau_{2k}(f)$ is in the same coset of $\text{im}(p_k)$ as $-\tau_k(W)$ and therefore $\tau_{2k}(f) \in \tau_{2k}(\Lambda)$ as desired.

Conversely, let $W \in \cT(\Lambda)$. Given any generating function of tube type $f:M \times \bR^{4N_1} \to \bR$ for $\Lambda$ with $W=T(f)$, Proposition \ref{prop: effect twisted stab} implies that for any fibered function of rigid tube type $q:M \times \bR^{4N_2} \to \bR$ the torsion class $\tau_{2k}(W)$ is equal to the torsion class $\tau_{2k}(T(f \oplus q))$  up to an element of $\text{im}(p_k)$. If we take $q$ so that $f \oplus q$ is frameable as in Proposition \ref{prop: framing for special Leg} then $\tau_{2k}(T(f \oplus q))$ and $\tau_{2k}(f \oplus q)$ are equal (the extra term in Theorem \ref{thm: framing} from the stable bundle vanishes since $\gamma_{f \oplus q}$ is trivial). Finally $\tau_{2k}(f)=\tau_{2k}( f \oplus q)$ because the fiberwise Morse theory is unchanged by stabilization (and the parity of indices is unchanged). Hence $\tau_{2k}(W)$ and $\tau_{2k}(f)$ are in the same coset of $\text{im}(p_k)$ and we are done. % up to sign, which is up to a sign equal to $\tau_{2k}(f)$. 
\end{proof}

\end{document}
